% Generated draft: requires mathematical and source review.
\documentclass[11pt]{article}
\usepackage[margin=25mm]{geometry}
\usepackage{fontspec}
\setmainfont{Latin Modern Roman}
\usepackage{amsmath,amssymb,mathtools,unicode-math}
\setmathfont{Latin Modern Math}
\usepackage{enumitem,needspace,xurl,hyperref}
\setlength{\parindent}{0pt}
\setlength{\parskip}{5pt}
\setlength{\emergencystretch}{3em}
\newcommand{\N}{\mathbb{N}}
\newcommand{\Z}{\mathbb{Z}}
\newcommand{\Q}{\mathbb{Q}}
\newcommand{\R}{\mathbb{R}}
\newcommand{\C}{\mathbb{C}}
\newcommand{\F}{\mathbb{F}}
\newcommand{\A}{\mathbb{A}}
\newcommand{\PP}{\mathbb{P}}
\newcommand{\E}{\mathbb{E}}
\newcommand{\eps}{\varepsilon}
\newcommand{\dd}{\,\mathrm{d}}
\begin{document}
\section*{UnsolvedMath ID: 20002335}
\textbf{Source problem code:} catalogue entry
\subsection{Definitions and mathematical statement}
Problem 26 asks when types in continuous logic are determined by the definable sets containing them, particularly whether \(\omega\)-stability suffices. Its background heading is empty. The distinction between a zero set and a distance-definable set is essential.

Let \(T\) be a complete continuous theory, \(A\) a parameter set, and \(S_{\bar x}(A)\) a finite-arity type space. Equip it with the logic topology and its canonical type metric
\[
 d(p,q)=\inf\{d(\bar a,\bar b):\bar a\models p,\ \bar b\models q\},
\]
using the prescribed tuple metric in a common sufficiently saturated model. An \(A\)-definable set here means a closed set whose distance predicate is \(A\)-definable; mere expression as a zero set is weaker. The type space is dictionaric when, for every logic neighborhood \(U\) of \(p\), there is such a set \(D\) with
\[
 p\in\operatorname{int}_{\mathrm{logic}}D\subseteq D\subseteq U.
\]
The aid attributes to Hanson the theorem that the relevant finite-arity type spaces of an \(\omega\)-stable continuous theory are dictionaric over every parameter set. Thus \(\omega\)-stability, in its continuous metric sense, gives a positive answer to the stated special question. This is not the classical assertion that there are literally countably many types.

The candidate representation applies to every dictionaric \(S_{\bar x}(A)\). For a nonempty \(A\)-definable \(D\), let
\[
 d(p,D)=\inf_{q\in D}d(p,q).
\]
This is the value of its definable distance predicate at \(p\). The aid reports
\[
 d(p,q)=\sup_D|d(p,D)-d(q,D)|,
\]
where \(D\) ranges over these definable sets. More precisely, for every \(0\le r<d(p,q)\), one can choose \(D\) with
\[
 p\in\operatorname{int}_{\mathrm{logic}}D,
 \qquad d(q,D)>r.
\]
Consequently the coordinate map
\[
 p\longmapsto(d(p,D))_D
\]
is isometric for the supremum coordinate metric and a homeomorphism onto its image for the logic topology and product topology. In the usual bounded-by-one convention the target is a product of \([0,1]\); another fixed metric bound changes that cube accordingly. The logic topology and metric topology need not coincide, so the two embedding assertions are kept separate.

In particular distinct types differ in membership in some distance-definable set. Arbitrary zero sets already separate types in every theory; that observation alone would not answer the problem. The aid supplies this sufficient condition and representation, not a classification of every theory with enough definable sets.
\subsection{Short English statement}
The aid reports that omega-stability gives enough distance-definable neighborhoods to distinguish types. In the resulting dictionaric spaces, all definable distance predicates reconstruct both the logic topology and the exact type metric. This is stronger than ordinary separation by arbitrary zero sets, which was never the difficult part.
\subsection{Sources}
\begin{itemize}
\item Local UnsolvedMath registry, permanent ID 20002335, exact source code catalogue entry; canonical \texttt{aim-logic-notes.json}, \texttt{notes[110]}, tag \texttt{problem}; AIM Workshop Problem Lists, Logic, \emph{Model theory of metric structures}, Problem 26. Created and updated 14 August 2026.
\item Source URL \url{https://aimath.org/WWN/continuouslogic/continuouslogic.pdf}, page 5 according to supplied verification. Its Background heading has no supplied content. The PDF was not accessed.
\item Hanson is named for the dictionaricness theorem, without a full title, date, theorem number, identifier, or separate URL.
\item Attempt 1: \texttt{partial\_result}, problem status at run \texttt{partially\_solved}, rigor label \texttt{proved}, novelty confidence \texttt{low}; classification \texttt{PARTIAL-PROGRESS}, registry status \texttt{partially\_solved}. Candidate: definable-distance norming/embedding theorem for dictionaric type spaces. These are dated supplied assessments, not independent certification.
\item Ulam AI / UnsolvedMath Contributors; CC BY 4.0. Machine-generated research aid, not peer reviewed; independently verify before citation.
\end{itemize}
\end{document}
